\label{sec:experiments_f}

\section{Experiments}
%\subsection{Datasets}
\label{sec:dataset_f}

% NOT USED . SEE EXPERIMENTS
Since there are no existing public datasets for contextual \pretty{Ic} and \pretty{Sl} task, we use two in-house datasets for evaluation - \textbf{booking} dataset,\footnote{We plan to release this dataset} which is a variation on \pretty{Dstc}-2 dataset \cite{williams2014dialog}, and \textbf{cable} dataset, a synthetically created multi-goal conversational dataset. 
 %Table \ref{data_table} shows statistics of these four datasets.
The \emph{Booking} dataset contains 9351 training utterances and 6727 test utterances, with 19 intents and 5 slot types. \emph{Cable} dataset comprises 1856 training utterances, 1814 validation utterances and 1836 test utterances, with 21 intents and 26 slot types.\footnote{Detailed Data statistics provided in supplementary material}
In addition, we also evaluate the model on non-contextual \pretty{Ic-Sl} public datasets - \pretty{Atis} \cite{hemphill90} and \pretty{Snips} \cite{coucke2018snips}.


%\subsection{Datasets}
%\label{sec:dataset_f}

\paragraph{Datasets:} Since there are no existing public datasets for contextual \pretty{IC} and \pretty{SL} task, we use two in-house datasets for evaluation - \textbf{Booking} dataset,\footnote{Dataset will be released to the public} which is a variation of \pretty{DSTC}-2 dataset \cite{williams2014dialog} with intent, slot and dialog act annotations, and \textbf{Cable} dataset, a synthetically created conversational dataset. 
 %Table \ref{data_table} shows statistics of these four datasets.
The Booking dataset contains 9,351 training utterances (2,200 conversations) and 6,727 test utterances, with 19 intents and 5 slot types. Cable dataset comprises 1856 training utterances, 1,814 validation utterances and 1,836 test utterances, with 21 intents and 26 slot types.\footnote{Detailed data stats provided in Appendix B}
In addition, we also evaluate the model on non-contextual IC-SL public datasets - \pretty{ATIS} \cite{hemphill90} and \pretty{Snips} \cite{coucke2018snips}.

\paragraph{Experimental setup:} 
To emphasize the importance of contextual signals in modeling, we first devise a non-contextual baseline of our \pretty{Casa-nlu} model, \pretty{Nc-nlu}. It is similar to \pretty{Casa-nlu} in terms of model architecture except no contextual signals are used. 
%the context fusion component is absent. 
Since neither datasets nor implementation details were shared in previous works on contextual NLU \cite{Bhargava2013EasyCI, Shi2015ContextualSL}, we also implement another baseline, \pretty{Cgru-nlu} that uses \pretty{GRU} \cite{cho2014learning} instead of self-attention for temporal context fusion. For fair comparison to existing non-contextual baselines, no pre-trained embeddings are used in any of the experiments though the model design can easily benefit from pre-training.
% Since neither datasets nor implementation details were shared in previous works on contextual NLU \cite{Bhargava2013EasyCI, Shi2015ContextualSL}, we first implement a baseline, \pretty{Cgru-nlu} that uses \pretty{GRU} \cite{cho2014learning} instead of self-attention for temporal context fusion.
% Moreover, to establish the strength of our chosen self-attention framework for this task, we further devise a non contextual baseline of our \pretty{Casa-nlu} model, \pretty{Nc-nlu}. It is similar to \pretty{Casa-nlu} in terms of model architecture except the context fusion component is missing. 

% \paragraph{Implementation Details:}
% We use hidden layer size of 56 with dropout probability of 0.3. Context history window $K$ was varied from 1 to 5 and the optimal value of 3 was selected.
% %,  Hyper-parameter optimization gave value of 3 for context history window $K$. 
% We use embedding layer size of 56 as well, training the embeddings from scratch. Adam \cite{kingma2014adam} algorithm with initial learning rate of 0.01 gave the optimal performance. Concatenation window size $P$ of 3 is used. $\alpha$ and $\beta$ in loss objective are set to 0.9. Early stopping is used with patience of 10 and threshold of 0.5. Each model is trained for 3 seeds and scores averaged across the seeds are reported.